\section{Podsumowanie projektu}

\subsection{Co zostało osiągnięte}
W ramach MVP dostarczono działający system \textit{end-to-end}, który:
\begin{itemize}
  \item zbiera zachowania użytkownika w aplikacji webowej,
  \item buduje cechy behawioralne z krótkich okien czasowych,
  \item klasyfikuje sesje \texttt{human} vs \texttt{bot},
  \item generuje decyzję enforcement na podstawie \texttt{risk\_score},
  \item umożliwia trenowanie i ewaluację modelu na danych lokalnych.
\end{itemize}

\subsection{Wnioski praktyczne}
\begin{itemize}
  \item Podejście behawioralne daje dobrą wykrywalność nawet przy botach imitujących człowieka.
  \item Interpretowalne drzewo decyzyjne ułatwia prezentację i audyt decyzji modelu.
  \item Największy wpływ nadal mają cechy dynamiki ruchu (zwłaszcza \texttt{acc\_max}),
  więc doskonalenie botów powinno iść w stronę realistycznego modelowania nagłych zmian tempa.
\end{itemize}

\subsection{Rekomendowane kolejne kroki}
\begin{enumerate}
  \item Rozszerzyć zbiór o nowe rodziny botów i sesje \textit{unseen}.
  \item Dodać walidację czasową (time-based split) oraz monitoring driftu cech.
  \item Uzupełnić policy engine o reputację źródeł i denylistę.
  \item Rozważyć model sekwencyjny (np. Transformer/RNN) jako warstwę drugą.
  \item Dodać dashboard operacyjny do przeglądu sesji i interpretacji predykcji.
\end{enumerate}
